%%%%%%%%% ABSTRACT
% Review note (claims vs. tables, not yet resolved): "two ways" (the paper has three remedies),
% "at no cost to the corruption gains" (the probe switch costs 0.4--7.1 corrupted points, Tab.~3),
% and "in every model" (the probe AUROC is reported for Qwen3-VL-4B only, Tab.~5).
\begin{abstract}
Vision--language models (VLMs) are deployed precisely where images are hardest to
read: assistive photography for blind and low-vision users, perception through rain,
fog and motion blur, and low-quality clinical or scientific scans. In these settings
the safe response to an uninterpretable input is not a confident guess but an
admission that the image cannot support an answer. We show that the obvious way to
prepare a VLM for degraded inputs---fine-tuning it for corruption robustness---silently
removes this behavior. Using two standard adaptation recipes, on-policy reverse-KL
distillation (RKL) and Group Relative Policy Optimization (GRPO), we find that
corruption-robust training improves clean and corrupted accuracy and transfers
zero-shot to corruption types never seen in training, yet drives both models to
answer uninterpretable and image-free inputs far more often than the base model.
The regression is invisible to accuracy: the most accurate model we train is the
worst at recognizing when it cannot see. We then explain the failure. A linear probe
recovers answerability from frozen hidden states at $\mathrm{AUROC}\approx0.93$ in every
model, so the ability is suppressed, not erased; and a reasoning-transplant test shows
the two recipes break for opposite reasons---RKL cannot reason its way to
``unanswerable,'' whereas GRPO reasons correctly but has learned an answer-level
shortcut that commits regardless. Finally we recover abstention two ways---adding
$10\%$ unanswerable examples to training, and inference-time activation steering toward
the base model---each at no cost to the corruption gains, with the training objective
deciding which works better. The results argue that corruption robustness and
calibrated abstention must be measured, and optimized, together.
\end{abstract}
